\documentclass[12pt,a4paper]{article}

\usepackage[utf8]{inputenc}
\usepackage[T1]{fontenc}
\usepackage[french]{babel}
\usepackage{amsmath,amssymb,amsfonts}
\usepackage{geometry}
\usepackage{enumitem}
\usepackage{xcolor}
\usepackage{hyperref}

\geometry{margin=2.5cm}

\newcounter{exercice}
\newenvironment{exercice}[1][Mathématiques]{%
    \refstepcounter{exercice}%
    \par\medskip\noindent
    \textbf{\large Exercice \theexercice\ -- #1}%
    \par\smallskip
}{%
    \par\medskip
}

\title{Annales de Mathématiques}
\author{Annales Probatoire}
\date{\today}

\begin{document}

\maketitle
\tableofcontents
\newpage

% Exercice: Exercice 0
% Sujet: Algèbre
% Généré automatiquement

\begin{exercice}[Algèbre]
\label{ex:1}

0,7^\{5\} pt\\
est asymptote à la courbe de\\
Justifie que la droile\\
2 a) Déterminer l'$expression / ' (x) de la fonction dérivée f$\\
0,5 pt\\
de\\
E;+æ[\\
\item b) Eludier le signe de f' (x) pour x €
0,5 pt\\
Dresser le tableau des variations de\\
0,7^\{5\} pt\\
sur\\
4^\{0\}\\
\item d) Construire la courbe de
+œ| et la compléter pour l'avoir sur R -\\
sur\\
pt\\

\end{exercice}

% --- Fin Exercice 0 ---


\newpage

% Exercice: Exercice 1
% Sujet: Analyse
% Généré automatiquement

\begin{exercice}[Analyse]
\label{ex:2}

1.a) Calculer la $limite de$\\
en +0\\
0,2^\{5\} pt\\
\item b) Calculer f'(r) où f' est la fonction dérivée de f .
0,5 pt\\
2.a) Dresser le tableau de variations de\\
[0; +œ[.\\
0,5 pt\\
su\\
\item b) Construire (C,).
0,5 pt\\
$Uo = 1 et pour$\\
\item 3. Soient (U,) et (\%;) les suites numériques définies respeclivement
par\\
3un\\
tout n € IN ,\\
Un+1\\
$V = 1 +$\\
3+4Un\\
5^\{0\}n+3\\
Montrer que pour tout entier naturel n, V + 4 =\\
0,2^\{5\} pt\\
est une suite arithmé$tique de raison 4 et de premier terme Vo = 4$.0,7^{5} pt\\
4.a) Montrer que\\
\item b) $Exprimer Vn en fonction de n pour tout entier naturel$ n.
0,2^\{5\} pt\\
En déduire Un en fonction de n-\\
0,5 pt\\

\end{exercice}

% --- Fin Exercice 1 ---


\newpage

% Exercice: Exercice 11
% Sujet: Algèbre
% Généré automatiquement

\begin{exercice}[Algèbre]
\label{ex:3}

1^\{1\} points\\
PROBLEME\\
Le problème comporte deux parties Aet B.\\
\item A) La courbe
r), ci-contre est la représentation graphique dune fonction numérique\\
dans un repère orthonormé $(0,5,1)$\\
Par lecture graphique\\
Déterminer l'ensemble de définition Df de\\
1^\{7\} et en 1*\\
,5 pt\\
ainsules $limites en$\\
2- Préciser le sens de variation de\\
dans\\
Résoudre\\
les inéquations\\
\item i) f(x)< 0
f(x)>0.\\
1pt\\
Déterminer f(-1) ; f(0) ; f(-1) où f est la\\
0,7^\{5\} pt\\
fonction derivée de\\
0,7^\{5\} pt\\
5- Dresser\\
tableau de variation de f\\
1l) On suppose que pour tout réel\\
4^\{1\}\\
$f(x)=ax +b+$\\
où a, b et c sont 3 réels.\\
En utilisant la question\\
montrer que\\
réels a, b et\\
vérifient le système\\
2b + c\\
0,7^\{5\} pt\\
=-2\\
2- Choisir la lettre correspondant\\
la bonne\\
0,7^\{5\} pt\\
reponse\\
Le triplet (a\\
\item c) est egal a
\item e) $(-2,1,8)$ ; 1) $(-2,4,-2)$ ; 9) $(1,2,4)$ ; h) $(0,-1, 0)$

\end{exercice}

% --- Fin Exercice 11 ---


\newpage

% Exercice: Exercice 2
% Sujet: Logique
% Généré automatiquement

\begin{exercice}[Logique]
\label{ex:4}

Exercice 2 $(4.5 points)$ _\\
Après un contrôle les notes de mathématiques de 6^\{0\} élèves dune classe de 1èrD\\
ont été regroupées dans le tableau suivant\\
Notes\\
[0; [\\
[1 ;8[\\
[8,1^\{2\}[\\
[1^\{2\} ;1^\{6\}[\\
[1^\{6\} ;2^\{0\}[\\
Nombre délèves\\
Fréquences\\
0,3\\
Effectifs cumulés croissants\\
[2pts]\\
Recopier et compléter ce tableau .\\
Calculer le pourcentage des élèves ayant une note supérieure ou égale\\
1^\{2\}/2^\{0\}\\
[0.Spt]\\
Lépreuve de mathématiques du contrôle est constituée de\\
exercices et dun problème.\\
Lenseignant dispose dans sa banque d'épreuve de 1^\{8\} exercices (dont\\
9 en\\
en statistiques\\
trigonométrie) et 1^\{0\} problèmes différents.\\
équations\\
Combien d'épreuves différentes peut-\\
[Ipt]\\
composer\\
Donner le nombre dépreuves contenant exactement un exercice de statistique et un exercice\\
de trigonométrie\\
[lpt] :\\

\end{exercice}

% --- Fin Exercice 2 ---


\newpage

% Exercice: Exercice 2X
% Sujet: Algèbre
% Généré automatiquement

\begin{exercice}[Algèbre]
\label{ex:5}

2x^\{2\}\\
que\\
f(x) =\\
On suppose\\
Ipt\\
Montrer que la droite d'é$quation y = -x$\\
1est une asymptote\\
à (C;) en\\
et +\%\\
symétrie de (CA).\\
Montrer que le point @(-1,0) est un centre de\\
Ipt\\
Construire (C;) dans le repère orthonormé $(0, 1,7 )$\\
Ipt\\

\end{exercice}

% --- Fin Exercice 2X ---


\newpage

% Exercice: Exercice 4
% Sujet: Algèbre
% Généré automatiquement

\begin{exercice}[Algèbre]
\label{ex:6}

EXERCICE 4\\
$(4,2^{5} points)$\\
Le plan est muni dun repère orthonormé $(0 ,1 ,J)$. La\\
courbe représentative (Cr) ci-contre est celle dune\\
fonction f , numérique à variable réelle.\\
En $exploitant cette courbe; déterminer$\\
\item a) L'ensemble de définition de f
0,2^\{5\}pt\\
\item b) Les $limites de f en$
0,5pt\\
c et en\\
1-\\
\item c) Une équation de l'asymptote verticale à la
0,2^\{5\}pt\\
courbe\\
f(-2), f (0) et f' (0).\\
0,7^\{5\}pt\\
\item e) Le sens des variations de
~2[ 0,5pt\\
sur ] -0\\
\item 2. On suppose que la fonction f est définie pour
tout x de son ensemble de définition par\\
f(x)\\
ax + b +\\
1+1\\
\item a) En $exploitant la question 1.d) , montrer$
1pt\\
$que a = 1$,$ b = 0 et c$\\
=1\\
\item b) Montrer que la droite (D) d'é$quation y = x est asymptote oblique $à
(c;)\\
0,5pt\\
\item 3. Calculer f'(x) pour tout x de l'ensemble de définition de f
0,5pt\\

\end{exercice}

% --- Fin Exercice 4 ---


\newpage

% Exercice: Exercice inconnu
% Sujet: Algèbre
% Généré automatiquement

\begin{exercice}[Algèbre]
\label{ex:7}

f(x)\\
III\_\\
La fonction\\
de la variable réelle\\
est définie sur R par\\
(C) désigne la courbe représentative de\\
dans le plan rapporté\\
un repère orthonormé\\
Calculer les $limites de$\\
0 et en+0\\
0,5 pt\\
Calculer la dérivée de\\
0,5 pt\\
Donner le tableau de variation de\\
(C) aux points\\
Écrire les équations des $tangentes$\\
et B dabscisses\\
et -1\\
respectivement\\
Tracer (C) ainsi que les $tangentes aux points$\\
et B\\
Voici quatre affirmations concernant la fonction f:\\
\item i) f est une fonction paire;
est une fonction positive sur R\\
iii) pour tout nombre réel x, 0 < f(x) < 3\\
iV)\\
est une fonction impaire.\\
Recopier sur votre feuille le tableau ci-dessous et le completer en met$tant une croix$\\
dans la case correspondant\\
votre choix selon que l'affirmation est vraie ou fausse\\
Affirmation\\
Vraie\\
Fausse\\
Activer\\

\end{exercice}

% --- Fin Exercice inconnu ---


\newpage

% Exercice: Exercice inconnu
% Sujet: Algèbre
% Généré automatiquement

\begin{exercice}[Algèbre]
\label{ex:8}

Partie C : 2 points\\
lintervalle [-4\\
la fonction définie sur\\
Soit\\
et dont une représentation graphique (4^\{2\}) dans\\
un repère orthonormé est donnée ci-contre\\
Déterminer par lecture graphique f(2), f(-3)\\
0,7^\{5\} pt\\
et f(O).\\
Résoudre graphiquement l'équation f(x)\\
$puis f(x) = 3$.\\
0,5 pt\\
Reproduire la courbe (0^\{2\})\\
et représenter la\\
courbe (9^\{2\}') de la fonction g définie sur\\
[-4\\
If(x)\\
6] par g(x)\\
0,7^\{5\} pt\\

\end{exercice}

% --- Fin Exercice inconnu ---


\newpage

% Exercice: Exercice inconnu
% Sujet: Algèbre
% Généré automatiquement

\begin{exercice}[Algèbre]
\label{ex:9}

La fonction\\
de la variable réelle\\
est définie pour tout nombre réel\\
différent de 3 par\\
2x +\\
f(x)\\
(Ct) désigne la courbe de\\
dans un repère orthonormé du plan\\
unité de longueur sur les axes est égale\\
cm\\
sur IR\\
Etudier les variations de\\
\{3\} -\\
construire le tableau de variation de f.\\

\end{exercice}

% --- Fin Exercice inconnu ---


\newpage

% Exercice: Exercice inconnu
% Sujet: Géométrie
% Généré automatiquement

\begin{exercice}[Géométrie]
\label{ex:10}

$f(x) = x^{2} - 1 $;\\
Soit les trois fonctions numériques\\
g et h définies sur IR par\\
$g(x) = x^{2} +$\\
h(x)\\
X +\\
repère orthonormé les courbes Cf\\
Tracer dans le plan rapporté\\
Cg et Ch ,\\
respectives des fonctions\\
g et\\
1,5 pt\\
Résoudre graphiquement le système\\
0,5 pt\\
solution de ce système\\
Pour tout\\
montrer que l'on peut toujours construire un triangle\\
ABC dont les longueurs des cotés sont\\
$BC = f(x)$\\
AC\\
AB\\
g(x)\\
h(x)\\
\item c)  Déterminer
tel que ABC soit isocèle de sommet principal B.\\
la fonction numérique définie par\\
t(x)\\
f(Ixl). A l'aide de la courbe Cf\\
Soit\\
Donner le programme de construction de la courbe Ct de la fonction\\
Tracer la courbe Ct dans le même repère que Cf.\\
0,5 pt\\

\end{exercice}

% --- Fin Exercice inconnu ---


\newpage

% Exercice: Exercice inconnu
% Sujet: Algèbre
% Généré automatiquement

\begin{exercice}[Algèbre]
\label{ex:11}

la fonction définie sur IR\\
Soit\\
par\\
$f(x) = 1$\\
X - 1\\
Soit Cfla courbe représentative de\\
et (H) la courbe d'équation Y =\\
(0,2^\{3\})\\
un repère orthonormé\\
dans le plan rapporté\\
Déterminer par ses coordonnées le vecteur de translation\\
transforme (H) en Cf.\\
qui\\
0,5 pt\\
Construire (H) et Cf dans le même repère.\\
1,5 pt\\
Canc\\
Vair\\
Bhidiar\\
Vahahong\\
Arofcar\\
Fahlaai\\
Wafiafinr\\
farrfor\\

\end{exercice}

% --- Fin Exercice inconnu ---


\newpage

% Exercice: Exercice inconnu
% Sujet: Algèbre
% Généré automatiquement

\begin{exercice}[Algèbre]
\label{ex:12}

un\\
unn\\
7nnn\\
7^\{291\}\\
Lcun\\
Acun\\
Déterminer graphiquement\\
intersection de Cf et de l'axe des abscisses .\\
labscisse du point I\\
0,2^\{5\} pt\\
é$quation f(x) = 2^{2}$\\
\item b) les solutions de
0,5 pt\\

\end{exercice}

% --- Fin Exercice inconnu ---


\newpage

% Exercice: Exercice inconnu
% Sujet: Algèbre
% Généré automatiquement

\begin{exercice}[Algèbre]
\label{ex:13}

points\\
Partle B\\
Soit (P) l'arc de la parabole définie sur [2\\
4] par\\
$f(x) = ax$' + bx + 1^{7} avec\\
et b sachant que (P) passe par le point B et que (P)\\
Déterminer les réels\\
admet pour sommet le point S$(3, -1)$\\
Soit 9, la fonction numérique de la variable réelle définie sur [2\\
4] par\\
2x^\{2\}\\
g(x)\\
1Zx +1^\{7\}\\
Etudier les variations de 9 et dresser son tableau de variations\\
Déterminer les abscisses des points F^\{1\} et Fz, intersection de (P)\\
avec l'axe des abscisses\\
Donner des équations cartésiennes des $tangentes aux points F^{1} et Fz$\\
Partie C :\\
points\\
Soit (H) la branche de I'hyperbole définie sur [4\\
+0\\
$h(x) = -1 +$\\
par\\
X - 3\\
Etudier les $limites aux bornes de [4$\\
+œ[\\
Dresser le tableau de variations de h-\\
Montrer que la droite d'équation\\
=-1 est asymptote horizontale\\
la courbe H\\
$1)^{2} = 4$\\
Représenter dans le même graphique le cercle d'équation\\
"arc de la parabole (P) et la branche de I'hyperbole (H) définie sur [4\\
\% [ ~\\

\end{exercice}

% --- Fin Exercice inconnu ---


\newpage

% Exercice: Exercice inconnu
% Sujet: Algèbre
% Généré automatiquement

\begin{exercice}[Algèbre]
\label{ex:14}

Étudier le sens de variation de g\\
1,5 pt\\
Étudier les $limites de$\\
+0 ,\\
0 et 1\\
1 pt\\
Dresser le tableau de variation de 9.\\
0,5 pt\\
Montrer qu'il existe 3 réels a, b et\\
$g(x) = ax + b +$\\
tels que\\
0,7^\{5\} pt\\
X-1\\
En déduire que (Cg) possède deux asymptotes dont on donnera les\\
3)\\
équations cartésiennes dans le repère (0, 7\\
0,5 pt\\
Tracer (Cg) dans $(0, 7,})$\\
On prendra\\
cm comme unité sur les axes\\
1 pt\\
Montrer que le point I(1, 3)\\
est centre de symétrie de (Cg)\\

\end{exercice}

% --- Fin Exercice inconnu ---


\newpage

% Exercice: Exercice inconnu
% Sujet: Algèbre
% Généré automatiquement

\begin{exercice}[Algèbre]
\label{ex:15}

Partie A\\
Lecture graphique\\
La représentation graphique de la fonction\\
est donnée ci-dessus\\
Déterminer graphiquement le domaine de définition de g\\
Résoudre graphiquement\\
(x)\\
g(x)\\
g(x)\\

\end{exercice}

% --- Fin Exercice inconnu ---


\newpage

% Exercice: Exercice inconnu
% Sujet: Algèbre
% Généré automatiquement

\begin{exercice}[Algèbre]
\label{ex:16}

PartieB :\\
5,5 points\\
La courbe Cr ci-contre représente une\\
fonction numérique f dans le repère\\
orthonormé $(0,1, J)$\\
a- Déterminer\\
conjecture,\\
par\\
Yensemble Di de définition de f.0,5pt\\
les\\
b- Déterminer par conjecture\\
$limites de f aux bornes de Dr$\\
Ipt\\
Dresser le tableau de\\
variation\\
Ipt\\
de f.\\
cahesierne ue Ullalurl\\
des asymptotes à Cf\\
Recopier et compléter le\\
0,7^\{5\}pt\\
tableau suivant\\
Equation ou inéquation\\
flr) <1\\
$flx) =0$\\
f^\{6\})^\{21\}\\
Ensomble des\_solutions\\
S,=\\
5^\{2\}=\\

\end{exercice}

% --- Fin Exercice inconnu ---


\newpage

% Exercice: Exercice inconnu
% Sujet: Algèbre
% Généré automatiquement

\begin{exercice}[Algèbre]
\label{ex:17}

Problème :$(llpoints)$\\
problème comporte deux parties\\
et B indépendantes.\\
Partie A\\
$(5,5points)$\\
2x =\\
[1;\\
+ol par f (x)\\
on note (C) la courbe représentative\\
On considère une fonction\\
définie sur\\
X+2\\
dans un plan muni dun repère orthonormé dunité 2cm.\\

\end{exercice}

% --- Fin Exercice inconnu ---


\newpage

% Exercice: Exercice inconnu
% Sujet: Algèbre
% Généré automatiquement

\begin{exercice}[Algèbre]
\label{ex:18}

Vérifier que sur F'intervalle [~I; +ol,$ f(x) = 2$\\
[0.Spt]\\
X+2\\
Calculer la $limite de$\\
en +o et en déduire l'existence dune asymptote Dà (C).\\
[0.Spt]\\
(a)   Calculer f' (x) où f\\
est la fonction dérivée de f\\
[0.Spt]\\
Dresser le tableau de variations de f\\
[Ipt]\\
Déterminer une équation de la $tangente (T) à (C) au point dabscisse 3.$\\
[0.Spt]\\
puis\\
Tracer les droites\\
et (D)\\
la courbe (C).\\
[Zpts]\\
(C') la courbe représentative de g\\
f(x)\\
Indiquer la transformation qui permet\\
On pose g (x)\\
dobtenir (€')\\
partir de (C), puis traver (C')\\
[lpt]\\

\end{exercice}

% --- Fin Exercice inconnu ---


\newpage

% Exercice: Exercice inconnu
% Sujet: Algèbre
% Généré automatiquement

\begin{exercice}[Algèbre]
\label{ex:19}

Le plan est muni dun repère orthonormé $(0,7,J)$ (unité\\
cm sur les axes)\\
Partie I : $(8 points)$\\
2r?\\
5x\\
f(x)\\
On considère la fonction f de la variable numérique x définic\\
(C) sa courbe\\
Par\\
2(x +\\
rcprésentative dans le plan, et Dr son ensemble de définition.\\
Détermincr Df-\\
[0,5 pt]\\
$Df f(x) = ax+b+$\\
Déterminer les récls a, b et\\
tels quc pour tout x élément dc\\
[1,5 pts]\\
Justifier que\\
est dérivable pour tout élément de Dr et calculer f' (x) .\\
[1 pt]\\
Déterminer les $limites de$\\
aux bornes de Df .\\
[1 pt]\\
Montrer que la droite dé$quation y = x + 2 est asymptote oblique $=\\
(C).\\
[0,5 pt]\\
Dresser le tableau de variation de f\\
[0,7^\{5\} pt]\\
maths.org\\
FUAEKruBlealeiiispuieranli\\

\end{exercice}

% --- Fin Exercice inconnu ---


\newpage

% Exercice: Exercice inconnu
% Sujet: Algèbre
% Généré automatiquement

\begin{exercice}[Algèbre]
\label{ex:20}

Montrer que le point I(-1;2\\
est centre de symétrie de (C).\\
[0,7^\{5\} pt]\\
Ac\\
Détermincr une équation cartésienne de la $tangente (T) à (C) au point d abscisse 0$\\
[0,5 pt]\\
\item 9. Tracer (C) ct (T)
[1,5 pts]\\

\end{exercice}

% --- Fin Exercice inconnu ---


\newpage

% Exercice: Exercice inconnu
% Sujet: Géométrie
% Généré automatiquement

\begin{exercice}[Géométrie]
\label{ex:21}

Probleme: $(1^{1} pomnts)$\\
Partic A\\
Le plan est muni dun repère orthonormé $(0, 7,])$\\
On considère la fonction rationnelle\\
définie par\\
$f(x) =x+1 +$\\
(Cf) sa courbe représentative dans le plan\\
(a)   Etudier les variations de f et dresser son tableau de variations.\\
[3pts\\
Préciser les asymptotes à la courbe (Cf) de f .\\
[0,Spt]\\
Démontrer que le point I(0; 1) est le centre de symétrie de la courbe (Cf)\\
[lpt]\\
[1,Spt]\\
Construire la courbe (Cr).\\
dé$finie par g(x) = f(lxl) et (Cg) sa courbe$\\
Soit g la fonction numérique de la variable réelle\\
[0,Spt]\\
representative\\
Montrer que la fonction g est paire:\\
[0,5pt]\\
partir de la courbe (Cr).\\
Donner un programme de construction de la courbe (Cg)\\
(b)\\
[0,Spt]\\
Tracer alors (Cg\\
repère que (Cr)\\
[0,5pt]\\
dans le même\\
Partie B\\
On désigne par A(1,3), B(-1,3) et C(-1,$ =1) trois points dans le rep$ère $(0, 7, )$)-\\
Quelle est la nature du triangle ABC?\\
[0,Spt]\\
Calculer $cos ACB et cos BAC$\\
En déduire les valeurs approchées en degré de ACB et BAC.\\
[2pts\\
MA^\{2\}\\
$MB^{2} = 2^{0}$\\
[lpt]\\
Trouver Pensemble (E) des points M du plan tels que\\

\end{exercice}

% --- Fin Exercice inconnu ---


\newpage

% Exercice: Exercice inconnu
% Sujet: Géométrie
% Généré automatiquement

\begin{exercice}[Géométrie]
\label{ex:22}

rartie\\
Le plan est muni d'un repere or\\
thonormé (0,\\
j). Le tableau ci-\\
Partie\\
deu tableau\\
contre est une\\
de variation dune fonction paire\\
f(x)\\
de courbe représentative (C).\\

\end{exercice}

% --- Fin Exercice inconnu ---


\newpage

% Exercice: Exercice inconnu
% Sujet: Analyse
% Généré automatiquement

\begin{exercice}[Analyse]
\label{ex:23}

[0,5 pt]\\
Donner le domaine de définition Df de\\
[0,5 pt]\\
(a) Déterminer les $limites de$\\
aux bornes de son domaine de définition.\\
Quels sont Yasymptote et lélément de symétrie de (C) ?\\
[1 pt]\\
[0,5 pt]\\
Quel est le signe de\\
0;0]\\
SUI\\
[1 pt]\\
Recopier et compléter le tableau de variation ci-dessus\\
(C). $(Unité sur les axes 2cm)$.\\
Construire la courbe\\
[1,5 pt]\\

\end{exercice}

% --- Fin Exercice inconnu ---


\newpage

% Exercice: Exercice inconnu
% Sujet: Analyse
% Généré automatiquement

\begin{exercice}[Analyse]
\label{ex:24}

On considère la suite (wn) définie par,$ Wn+1 =b(c)$" +bn+a\\
Pour tout ne IN, on pose\\
\item B) | -
partie\\
A-I^\{1\})\\
$Un =b(c)$'\\
$et Vn= bntao$ù\\
b et\\
sont les reels dela\\
géométrique et (Va) une suite arithmétique\\
suite\\
est\\
Montrer que (un)\\
une\\
2^\{2\}n+1 +2n+1\\
0, Spt\\
Montrer aue pour tout entier naturel\\
Yr\_\\
$Expnmer Sn puIs $ ,$\\
toncton de\\

\end{exercice}

% --- Fin Exercice inconnu ---


\newpage

% Exercice: Exercice inconnu
% Sujet: Géométrie
% Généré automatiquement

\begin{exercice}[Géométrie]
\label{ex:25}

Partic B\\
Le plan est rapporté au repère orthonormé $(0,1.7)$\\
Soit f la fonction dé$finie sur D = R $= {1 } par f(x)\\
(C;) la courbe\\
On note\\
representative de f\\
1.a) Calculer les $limnites de$\\
aux bornes de D\\
(Cp).\\
0,2^\{5\} pt\\
\item b) En déduire une asymptote
la courbe\\
Dêterminer trois recls a.b et\\
de D.fx)\\
tels que pour tout\\
b +\\
0,7^\{5\} pt\\
la courbe (Cà\\
\item d) Montrer que la droite (1) :$ v = 1$
0,2^\{5\} pt\\
Csl asvmnptote\\
\item 2. a) Déterminer la fonction dérivéc de
et dresser le tableau de variation de \}\\
0,7^\{5\} pt\\
\item b) Existe-t-il des points de (C)) oü la $tangente à (Cp) est parallèle$
la droite (4) ? Justifier\\
Votre réponse\\
0,5 pt\\
Tracer la courbe (Ca)\\
0,7^\{5\} pt\\
inspirer de la courbe (C;). Soit m un paramètre réel On\\
Pour ce\\
4u/\\
Suil, On pourra\\
considere\\
équation (En)\\
0 et on note An son discriminant .\\
mr\\
4,0)\\
Calculer 4\\
0,2^\{5\} pt\\
\item h) Pour quelles valeurs de m Féquation (En) admet-elle une solution unique
0.5 pt\\
\item c) Pour quelles valeurs de m^\{1\}
équation (Em) admet-elle deux solutions distinctes\\
0,5 pt\\
On suppose que Féquation (En) admet deux solutions distinctes X^\{1\} Ct 8^\{2\}.\\
$Exprimer,$\\
produit\\
en fonction de m. la somme Sm et le\\
Pm de ces solutions ,\\
0,5 pt\\
\item b) Peut-on trouver ml tel
que\\
les solutions X^\{1\} et\\
soient toutes négatives\\
0,5 pt\\
Pour\\
quelles valeurs de m les solutions 1p et X^\{2\} sont-elles de signes contraires\\
0,5 pt\\
Pour\_quelles\_Valeurs\_de m les solutions XLet Xz sont-clles\_toutes\_positives\\
0.5 pt\\

\end{exercice}

% --- Fin Exercice inconnu ---


\newpage

% Exercice: Exercice inconnu
% Sujet: Algèbre
% Généré automatiquement

\begin{exercice}[Algèbre]
\label{ex:26}

EXERCICE\\
Vuiis\\
ZY+4\\
et on désigne par\\
Lm;vl + lla fonction\\
J(x) =\\
On detinit sur\\
par\\
2^\{173\}\\
courbe dans un repère orthonorme $(0,i,1)$. Unité sur les axes\\
Cmnl\\
symétrie de Cr:\\
0,5 pt\\
0(\\
est un centre de\\
\item a) Demontrer que le point
0,5 pt\\
$limite de$\\
en +œ\\
\item c) Demontrer que la droite

\end{exercice}

% --- Fin Exercice inconnu ---


\newpage

% Exercice: Exercice inconnu
% Sujet: Algèbre
% Généré automatiquement

\begin{exercice}[Algèbre]
\label{ex:27}

Exercico\\
$(5 points)$\\
Chacuno dos c$lna quostions ci-après a qualre propositlons de ropo$nse. Écrire le numoro\\
de la question suiM de la lettre indiquant Ia réponse Juste. Chaque question est notée\\
pL\\
sur\\
Soit f la fonction dofinie sur R - (~$2) par f(x) = x -1$\\
Pour tout x de\\
$R = {$~2)que vaut f'(r) ?\\
\item 6) 4
4^\{4\}\\
\item c) 4"4
2)\\
G+z\\
(+2)\\
On considère léquation (E):\\
2)\\
$sinx + VBcosx = 1$. Quelle est l'équatlon équivalente\\
0 l'équation (E) ?\\
\item b) sin(*-\#$) =-$
\item a) $sin (x$
'(+4)\\
(-4) \}\\
\item H) =
\item d) sin
\item c) sln
$Soit / la fonction deflnie sur R par fGz) = x$'. On considoro la fonctlon 9 définie sur\\
3)\\
$Rpar g(r) = (r + 2)$? + 1. Dans lo plan muni dun repore, la courbe de la fonction 0\\
'Timage de la courbe de la loncton f par la translation de vecteur 0. Quel est Ie\\
esi\\
couple de coordonnges du vectour ù ?\\
\item a) (2;1) ;
\item b) (2;-1)
\item c) (~2;1) ;
\item d) $(~2;-1)$.
Soit h la fonction définle sur R - \{1) par h(x) =\\
Quello est la $limlte €e h(x)$\\
à gaucho ?\\
lorsque x lend vors\\
\item 0) 3
\item 6) 3
\item d) +\%
\item c) ~0
Tous les parucipanls à uno conforonco\\
échangont 3^{25} polgnoos do ma$ln. Deux$\\
porsonnes quelconques n'échangent qu'uno soule polgnoo de main. Quel ost le\\
nombre de porsonnes ayant asslst^\{6\} € la conlorence\\
Activer\\
\item b) 2^\{6\}
\item 0) 2^\{5\} ;
\item c) 2^\{7\}
\item d) 2^\{4\},
Accedez\\

\end{exercice}

% --- Fin Exercice inconnu ---


\newpage

% Exercice: Exercice inconnu
% Sujet: Algèbre
% Généré automatiquement

\begin{exercice}[Algèbre]
\label{ex:28}

Exercice\\
($(5 points)$\\
Chacuno dos c$lna quostions ci-après a qualre propositlons do ropo$nse. Écrire le numoro\\
de la lettre indiquant Ia réponse Juste. Chaque question est notée\\
sUIM\\
de la question\\
PL.\\
sur\\
Soit f la fonction dofinle sur R - (~$2) par f(x) = x -1$\\
Pour tout x d^\{0\}\\
R - (~2)que vaut f'(x)\\
\item 9) 4^\{4\}
9=\\
6r\\
6+2)\\
On considère réquation (E):\\
2)\\
$sinx + VBcosx = 1$. Quelle est l'équatlon équivalente\\
0 l'équation (E) ?\\
(-4)-\}\\
\item b) $sin (x$
3)\\
(+;) =\\
\item a) sin
\item c) sln
\item d) sln
Soit / la fonction def$lnie sur R par /G)$\\
On considoro la fonctlon 9 définie sur\\
$Rpar g(r) = (r + 2)$? + 1. Dans lo plan muni dun repore, la courbo de la fonction 0\\
est $lImaoe de la courbe de la lonction$\\
Dar la translalion da vartenr 7\\
Oual est le\\

\end{exercice}

% --- Fin Exercice inconnu ---


\newpage

% Exercice: Exercice inconnu
% Sujet: Algèbre
% Généré automatiquement

\begin{exercice}[Algèbre]
\label{ex:29}

Partie A\\
Soit f la fonction dé$finie par f(x)=-x+2 -$\\
et (C) sa courbe représentative\\
$X = 2$\\
dans un repère orthonormé $(0 ,7,1)$. (unités sur les axes\\
cm)\\
\item 1) Déterminer l'ensemble de définition de f
0,Spt\\
2)a) Calculer\\
$limites de f en$\\
2 et 2*\\
1pt\\
les\\
\item b) Justifier que la droite d'é$quation x = 2 est asymptote verticale $à (C).
0,5pt\\
\item 3) Calculer f '(x) et justifier que f '$(r) 2^{0} =xe[1$
2[UJ^\{2\}\\
3]\\
Ipt\\
4)a) Justifier que la droite d'équation y =\\
0,Spt\\
2 est une asymptote oblique à (C).\\
\item b)   Tracer (C) et ses asymptotes dans le repère orthonormé $(0,i,j)$.
2pts\\
Soit m un nombre réel quelconque\\
Déterminer suivant les valeurs de m\\
1,5pt\\
le nombre et le signe des solutions de l'équation f(x)\\

\end{exercice}

% --- Fin Exercice inconnu ---


\newpage

% Exercice: Exercice inconnu
% Sujet: Algèbre
% Généré automatiquement

\begin{exercice}[Algèbre]
\label{ex:30}

FARDEA\\
opoins\\
dune variable réelle x dont le tableau de\\
On considère la fonction numérique\\
variation est donné ci-dessous\\
+\%\\
f(x)\\

\end{exercice}

% --- Fin Exercice inconnu ---


\newpage

% Exercice: Exercice inconnu
% Sujet: Algèbre
% Généré automatiquement

\begin{exercice}[Algèbre]
\label{ex:31}

On désigne par (C;) sa courbe représentative dans un repère orthonormé $(0, 7 ,)$):\\
0,2^\{5\}pt\\
Déterminer lensemble de définition de\\
0,2^\{5\}pt\\
Déterminer le signe de f(x) pour x € ]-1, +œ[\\
0,2^\{5\}pt\\
Déterminer une asymptote à (Ci).\\
0,5pt\\
Déterminer une équation de la $tangente à (Ci) au point dabscisse$\\
On suppose que f(x)\\
ax + b +\\
système ci-dessous\\
0,7^\{5\}pt\\
Montrer que $(a,b,c)$ est solution du\\
$y + 2 = -2$\\
d'inconnue $(x,y,z)$\\
$X = z $= 0\\
$y + z = $=2\\
2x\\
b)\_En\_déduire a. b et c.\\
Ipt\\

\end{exercice}

% --- Fin Exercice inconnu ---


\newpage

% Exercice: Exercice inconnu
% Sujet: Géométrie
% Généré automatiquement

\begin{exercice}[Géométrie]
\label{ex:32}

Partie B\\
Représentation graphique de\\
X^\{2\} + X | 2\\
On suppose que f(x)\\
X | 2\\
tels que f(x)\\
Déterminer trois réels a, b et\\
0,7^\{5\} pt\\
b +\\
4 2\\
(D) d'é$quation Y = x -$\\
est asymptote à la courbe (Cf) de f.\\
que\\
Montrer\\
la droite\\
0,5 pt\\
Vinfini.\\
Calculer les $limites de$\\
au point -2\\
1 pt\\
Étudier le sens de varations et dresser le tableau de vanations de f\\
1,5 pt\\
Écrire une équation de la $tangente (T) à$\\
au point d'abscisse -1\\
la courbe de\\
0,5 pt\\
les droites (T), (D) et la partie de la\\
Tracer dans\\
repère orthonormé du plan\\
pts\\
courbe (Cf) correspondant\\
lintervalle ]-2\\
+œ[.\\
$f(2) = 2$\\
Vérifier que f(-1)\\
résoudre graphiquement dans l'intervalle ]-2\\
puis\\
+œ[:\\
0,7^\{5\} pt\\
f(x)\\

\end{exercice}

% --- Fin Exercice inconnu ---


\newpage

% Exercice: Exercice inconnu
% Sujet: Algèbre
% Généré automatiquement

\begin{exercice}[Algèbre]
\label{ex:33}

Partie B\\
On considère l'équation (E) : - x?\\
$x(4 = m) - 5 + 2m $= 0 où x est l'inconnue et m un\\
paramètre réel.\\
\item 1) Justifier que 2 ne peut pas être solution de cette équation quelque soit m.
0,5pt\\
\item 2) Montrer que pour tout réel x distinct de 2,
$2m = 0 $~$f(x) = m$.\\
0,7^\{5\}pt\\
x(4\\
\item $m) = 5$
3)a) Calculer le discriminant de (E) en fonction de m.\\
0,7^\{5\}pt\\
léquation (E) n'admet pas de solution si et seulement si\\
que\\
\item b) En déduire
me]-2\\
1pt\\
2[\\
\item c) On suppose que me]-2
2[. Déterminer par calcul, les valeurs de m pour\\
lesquels l'équation (E) admet deux solutions de signes contraires\\
Ipt\\

\end{exercice}

% --- Fin Exercice inconnu ---


\newpage

% Exercice: Exercice inconnu
% Sujet: Analyse
% Généré automatiquement

\begin{exercice}[Analyse]
\label{ex:34}

Exercice\\
points\\
Lableel slivanl donne\\
reparliuion des candidel\\
(lalheniauale sel^\{0\}\\
Conlours\\
Ieur> [oles\\
6, 1o[\\
[oar\\
Lladg\\
Se[\\
5^\{16\} 1^\{8\}[\\
5^\{1830\}\\
Eeehrs\\
Shno\\
YkEuis (ilils\\
Cicissng\\
YkEuis (ilils\\
Recopier\\
Complelen\\
Lableau Pef-\\
Igrie Uesrejuenles\\
eetils cumules Ctoissenl\\
ligne des elfeclifs cumule :\\
0,7^\{5\} pE\\
Jegdissanls\\
Reponure Sar Yrai\\
ue Cuciddall onl\\
gienolegudeneugy\\
1^\{20\} candidals\\
0,2^\{5\} 2^\{2\}\\
une ndle Sucerilure\\
Elk\\
0,2^\{5\}\\
5^\{0\}w^\{6\} des cendidals unl\\
Une noleinierieure\\
Cortruicedar\\
MeMeeote\\
Folygdres des eitcjig Cumules joissents\\
Lecdissanls\\
(On prendre\\
ersCisses\\
uniles\\
Jiues\\
$Sin Vui$\\
0aonneee)\\
0,5 pl\\
Deleritiirier\\
Llasse Modale\\
Lelle serie slelislique\\
Delefilinen\\
Jechiouernenl\\
(iiejiene\\
Celle Sene Slalisuidve\\

\end{exercice}

% --- Fin Exercice inconnu ---


\newpage

% Exercice: Exercice inconnu
% Sujet: Analyse
% Généré automatiquement

\begin{exercice}[Analyse]
\label{ex:35}

Exefcice\\
Doinls\\
Les Covfbes\\
demandees seront (racees Sur un Babier Mil$limetre:$\\
Lesler leur dutee de\\
On elledtue Ues Essais sur$ln echanlillcoi de 2^{20} lempes Eleclriques elin$\\
eue duree\\
$exprirnee$\\
Deureo\\
rEsuilals SDill (eoidupes Dar 4iasses Uanchluue euaie\\
heures UAn\\
Lableannvanl\\
[1 / : 1 b\\
(1,4 ; 1 9\\
Vuge^\{4\}s\\
1^\{3\}: La^\{4\}i\\
1io : 2n\\
rilliers\\
Eleslils\\
freouences\\
Cumulees\\
Ldissanle:\\
Suplose Ole\\
[eparliucr\\
unotrhe\\
Mgileneue\\
Chauve €a^\{54\}\\
TrazerT'histourarniie\\
Cel Senie\\
Recboier\\
Lunlcieler\\
LablealG-dessus\\
Calculer\\
valeur apprulhet\\
tnediane\\
Uelaur\\
LeuesAcie\\
Atr\\
Polyuun\\
desireaediees Cumnule^\{4\}s Coisrenies\\
Celculer\\
[ronlalures Vordre\\
eldelecart(ype\\
serie\\
Mleyenne\\
lenres dort\\
duree\\
Vuurcenlaul\\
aulreal\\
Jarnpes eleciques\\
tgne duisgenee Drovenpdk degiaulre\\
haorcanl\\
Yleaienieailresy\\
Mcyen\\
danree\\
ecani\\
0,5 pl\\
celui Ues Deux Iols Qui Yuus semnble elre ineilleur\\

\end{exercice}

% --- Fin Exercice inconnu ---


\newpage

% Exercice: Exercice inconnu
% Sujet: Algèbre
% Généré automatiquement

\begin{exercice}[Algèbre]
\label{ex:36}

Un comité de développement d'un village voudrait acheter les appareils suivants\\
une tronçonneuse et un groupe électrogène\_\\
une motopompe\\
Pour obtenir des fonds\\
B et C selon leurs revenus\\
repartit ses membres en trois groupes A\\
Le tableau ci-dessous\\
donne la contribution de chaque membre par appareil en fonction de son groupe\\
Contribution par membre\\
Appareils\\
Groupe\\
Groupe B\\
Groupe\\
Motopompe\\
4^\{500\}\\
7^\{500\}\\
1^\{2000\}\\
Tronçonneuse\\
7^\{000\}\\
1^\{0000\}\\
0^\{00\}\\
Groupe électrogène\\
3^\{500\}\\
4^\{500\}\\
6^\{000\}\\
Sachant que la motopompe, la tronçonneuse et le groupe électrogène coûtent\\
5^\{46\} 0^\{00\} F\\
respectivement\\
7^\{70\} 0^\{00\}\\
3^\{42\} 0^\{00\} F\\
Calculer le nombre de membres de chaque groupe\\
En déduire le nombre de membres de ce comité\\
0,5 pt\\

\end{exercice}

% --- Fin Exercice inconnu ---


\newpage

% Exercice: Exercice inconnu
% Sujet: Analyse
% Généré automatiquement

\begin{exercice}[Analyse]
\label{ex:37}

On procède à une nouvelle répartition par tranche d'âge des membres de ce\\
comité et on obtient le tableau suivant\\
âge\\
[3^\{0\};3^\{5\}[\\
[4^\{5\}\\
Tranche\\
[3^\{5\}\\
4^\{0\}[\\
[4^\{0\}\\
4^\{5\}\\
5O^\{1\}\\
[5^\{0\}\\
5^\{51\}\\
Effectifs\\
Donner les classes modales \_\\
0,5 pt\\
Un seul des couples ci-dessous représente la moyenne x et la valeur approchée\\
"'écart type\\
près\\
1^\{0\}\\
de cette série statistique \_\\
Ecrire la bonne réponse\\
sur votre feuille de composition\\
(4^\{2\}\\
7,0^\{5\})\\
$(5^{2} ; 8,2^{5})$\\
\item 4) (4^\{2\}
8,2\\
(3^\{9\}\\
7,0^\{5\})\\

\end{exercice}

% --- Fin Exercice inconnu ---


\newpage

% Exercice: Exercice inconnu
% Sujet: Mathématiques
% Généré automatiquement

\begin{exercice}[Mathématiques]
\label{ex:38}

Exercke\\
points\\
Les ncles des eleves jure Llasse De\\
reperlies suivanlel\\
pericutarle\\
Curs\\
dune Evalualion   Sunl [ediesenlees Dants\\
lableedci-dessous\\
Ko :\\
ENectls\\
epprochees\\
Jes(esulal\\
Ondonnet\\
Vael\\
Pr^\{2\}s Pei\\
EehS\\
oeleriigier\\
2i^\{2\}rseode\\
inesausuoue\\
Conslruire|hislouranune Des elleclils\\
1,Spt\\
Calculer\\
Mde\\
Mlediarie\\
Moyenne\\
Tecarl lyre\\
eul\\
Vagiegiee\\

\end{exercice}

% --- Fin Exercice inconnu ---


\newpage

% Exercice: Exercice inconnu
% Sujet: Mathématiques
% Généré automatiquement

\begin{exercice}[Mathématiques]
\label{ex:39}

Exercice\\
points\\
fegt^\{152\} dans\\
aqenl Ieceiseli\\
Orlarne\\
ecelles descoininecanl\\
\#xciiyes\\
tilliers\\
[ranes ela iekve\\
lackav Sunan\\
[2^\{5\} ; J^\{05\}\\
[2^\{0\} ; J^\{5\}[\\
[1^\{5\}   2^\{0\}\\
[2^\{0\} : 2^\{5\}\\
Ehecuiis W\\
Dateunineshses^\{4\}\\
tnd\\
itogramrire des e(lectib\\
Kansr~\\
Evaher\\
$lnambie$\\
coinlfrvrcants dant\\
iecelleiouualyie\\
es^\{1\} Coindiieenle 2^\{5\}\\
2t^\{40\} mrills \{ranes\\
Cacuk\\
akurmnoyenedes iecelle:\\
purnaliresde\\
2^\{23\} Coiniitieicanl:\\
dixyramririe des ellcuils\\
Consunife\\
diaarairurie des ellelib Clinlks Cioishnl\\
Clinles deciarsantssui\\
ityinieatachnue.\\
0,5 2l\\
Delenniiner Qtachnveiynt\\
irvdiane\\
shi^\{2\}.\\

\end{exercice}

% --- Fin Exercice inconnu ---


\newpage

% Exercice: Exercice inconnu
% Sujet: Probabilités
% Généré automatiquement

\begin{exercice}[Probabilités]
\label{ex:40}

On considère le tableau suivant :\\
Classes\\
[2^\{5\};3^\{0\}[\\
[1^\{5\};2^\{0\}[\\
[2^\{0\};2^\{5\}[\\
[3^\{0\},3^\{5\}[\\
[3^\{5\};1^\{0\}[\\
Elfectifs\\
Elfectifs\\
cumulés\\
croissants\\
Elfectifs\\
cumulés\\
décroissants\\
Recopier et compléter ce tableau.\\
[2 pts]\\
Construire sur un même graphique le diagramme des cffectifs cumulés croissants, ct cclui des\\
effectifs cumulés décroissants\\
[2 pts]\\
approchée de la médiane de cette séric statistique.\\
[1 pt]\\
En déduirc unc valeur\\

\end{exercice}

% --- Fin Exercice inconnu ---


\newpage

% Exercice: Exercice inconnu
% Sujet: Géométrie
% Généré automatiquement

\begin{exercice}[Géométrie]
\label{ex:41}

EaCnlilc\\
puIILS\\
Dans une classe de première D comp$tant 9^{0} élèves dont 6^{0} garçons$\\
une enquête\\
est menée sur la dis$tance hebdomadaire en km parcourue par chaque élève p$our se rendre\\
au Lycee\\
Le resultat est consigne dans le tableau complet cI-dessous\\
Dis$tances$\\
7;1^\{1\}\\
Effectifs\\

\end{exercice}

% --- Fin Exercice inconnu ---


\newpage

% Exercice: Exercice inconnu
% Sujet: Géométrie
% Généré automatiquement

\begin{exercice}[Géométrie]
\label{ex:42}

Ea^\{5\}l^\{9\}l^\{5\}\\
PuMllS\\
Une enquête révèle que 1^\{6\}\\
Une entreprise emploie 5^\{0\} personnes dont 2^\{0\} femmes\\
ont une assurance maladie.\\
employés de cette entreprise; dont 5 femmes\_\\
1,Spt\\
Recopier et compléter le tableau ci-dessous\\
Sans assurance\\
Avec assurance\\
Total\\
maladie\\
maladie\\
2^\{0\}\\
Femmes\\
Hommes\\
5^\{0\}\\
Total\\
1^\{6\}\\
Les employés de cette   entreprise décident de protester contre leurs conditions de\\
négocier.\\
Ils   décident alors\\
faire\\
de\\
L'employeur\\
se\\
travail.\\
de\\
leur\\
propose\\
négociations .\\
représenter\\
table\\
une délégation de\\
personnes\\
des\\
par\\
le nombre de différentes délégations\\
dans chacun des cas suivants\\
Déterminer ,\\
possibles qu'ils peuvent constituer\\
\item 1) Chaque délégation comporte au moins une femme sans assurance
0,7^\{5\}pt\\
maladie\\
\item 2) Chaque délégation comporte exactement une femme et exactement deux
0,7^\{5\}pt\\
personnes sans assurance maladie.\\
On regroupe maintenant les employés de cette entreprise suivant leurs salaires\\
journaliers en francs CFA\\
5ooo[\\
[4^\{000\}\\
7^\{000\}[\\
[7^\{000\}\\
T^\{00\}oo[\\
5^\{000\}\\
2^\{000\}\\
3^\{000\}[\\
[3^\{000\}\\
4^\{000\}[\\
Salaires\\
2^\{0\}\\
3^\{0\}\\
1^\{6\}\\
1^\{0\}\\
Fréquences\\
\item 1) Construire le diagramme des fréquences cumulées décroissantes de cette
1pttiver !\\
série\\
eoez a^\{1\}\\
\item 2) Calculer le salaire journalier médian des employés de cette entreprise.
0,5pt\\

\end{exercice}

% --- Fin Exercice inconnu ---


\newpage

% Exercice: Exercice inconnu
% Sujet: Analyse
% Généré automatiquement

\begin{exercice}[Analyse]
\label{ex:43}

$(5 points)$\\
Excrcice\\
Le tableau suivant donne la répartition des tailles $(en centimetres)$\\
des cleves dune classe de\\
[*re TI\\
[4^\{5\}. ISS[\\
Taille\\
1^\{153\}\\
6^\{01\}\\
7l^\{60\}.\\
1^\{6\}s^\{1\}\\
[l^\{63\}.\\
1^\{706\}\\
\item 1^\{170\}. 1^\{90\}[
Effectif\\
repere convenablement choisi. Thistogramme de cette série\\
Tracer; dans\\
1,5 pt\\
polygone des effectifs cumulés croissants de cette série. Pour cela, on prendra\\
2, a) [racer ]e\\
individus.\\
enl abscisse\\
cIl cl en ordonnce\\
1,2^\{5\} pt\\
Cin) Pour\\
cm pour\\
\item b) Déterminer eraphiquement et par calcul au centimetre près, la taille n^\{1\} médiane de ces
graphique.\\
élèves On laissera apparent les trais utilises pour la lecture\\
1,2^\{5\} pt\\
Calculer la moyenne m^\{71\} etla variance\\
de cette scrie\\

\end{exercice}

% --- Fin Exercice inconnu ---


\newpage

% Exercice: Exercice inconnu
% Sujet: Analyse
% Généré automatiquement

\begin{exercice}[Analyse]
\label{ex:44}

Eyorcico\\
points\\
les regroupe en classes a'arplltude\\
Les effectifs des électeurs des tranches dâges [2^\{0\} ,3^\{5\}[ et [3^\{5\};5^\{0\}[ sont\\
respectivement 2^\{1\} et 1^\{9\}.\\
Les classes modales de cette série sont [3^\{0\} , 3^\{5\}[ et [3^\{5\} , 4^\{0\}[ et ont pour effectifs\\
respectifs 1^\{0\}.\\
Les 3 électeurs de la classe [2^\{0\} , 2^\{5\}[ et les 2 électeurs de la classe[4^\{5\} , 5^\{0\}[ sont les\\
seules femmes de ce bureau de vote\_\\
1,5pts\\
Dresser le tableau des effectifs de cette série regroupée en classes\\
1,Spts\\
Construire le polygone des effectifs cumulés décroissants\\
(Icm pour\\
3 unités)\\
Ipt\\
Déterminer graphiquement et par calculs la médiane de cette série\_\\
Les 4 premiers vo$tants de ce bureau seront réunis pour former un$\\
comité de surveillance\\
0,Spt\\
\item a) Déterminer le nombre de comités possibles
0,5pt\\
\item b) Déterminer le nombre de comités compor$tant 2 femmes.$

\end{exercice}

% --- Fin Exercice inconnu ---


\newpage

% Exercice: Exercice inconnu
% Sujet: Mathématiques
% Généré automatiquement

\begin{exercice}[Mathématiques]
\label{ex:45}

dun certain nombre de lapins a donné le\\
poids\\
jardin; une' observation des\\
Dans un\\
résultat suiyant\\
Erecuis\\
sur 2\\
Page\\

\end{exercice}

% --- Fin Exercice inconnu ---


\newpage

% Exercice: Exercice inconnu
% Sujet: Analyse
% Généré automatiquement

\begin{exercice}[Analyse]
\label{ex:46}

\item 1) Déterminer le poids moyen de ces lapins.
$(0,7^{5} pt)$\\
encore appelé polygone des effectifs\\
\item 2) Construire la courbe cumulative décroissante
$(1,5 pt)$\\
cumulés décroissants.\\
$(0,7^{5} pt)$\\
\item 3) Déterminer la médiane de cette série statistique.

\end{exercice}

% --- Fin Exercice inconnu ---


\newpage

% Exercice: Exercice inconnu
% Sujet: Géométrie
% Généré automatiquement

\begin{exercice}[Géométrie]
\label{ex:47}

EXERCICEZ\\
points)\\
Le lableau ci-dessous regroupe les notes sur 1^\{00\}, obtenues\\
8^\{0\} candidats\\
un lest\\
ecnl\\
de selection .\\
Notes\\
[0;\#S[\\
[4^\{5\} ; 5o[\\
[5^\{0\} :6^\{0\}[\\
[7^\{5\} ;B^\{5\}L\\
uue^\{3\}\\
Effectifs\\
Cakuler\\
nole moyenne des candldats\\
test de selection\\
médiane de cette sene par interpolaton lineaire\\
cakculer\\
candidats dont trcis fcmmcs\\
ayant oblenu une note superieure ou egale\\
esSi\\
son; soumis\\
\{est oral\\
On decido dc commencer\\
groupe de trois candidats\\
slmulianemcol\\
cho 5^\{15\}\\
Délerminer\\
nombre de grcupes\\
peut former,\\
oossibles Queron\\
Japl\\
Determiner\\
possibles comprenant au moins\\
emmes\\
nomorcdogroupes\\
Pour\\
Les six candidats retenus\\
soni Invites\\
stasseoir autour dune table d^\{0\}\\
te le scrte quc deux\\
ient pas\\
Six cha ses\\
assisse^\{5\} cole\\
Ocisonncs\\
mame sexc\\
loules les femmes\\
Les femmes\\
a-colc\\
Chacun des trois hommes Ien^\{0\}\\
maun\\
saluentoasenbre\\
Dessiner\\
graphe pcrmct$tant de modeliser cette situation.$\\
\item b) Quel est
degre dc chaque sommet de ce\\
graphe\\
0,5pt\\

\end{exercice}

% --- Fin Exercice inconnu ---


\newpage

% Exercice: Exercice inconnu
% Sujet: Analyse
% Généré automatiquement

\begin{exercice}[Analyse]
\label{ex:48}

$(0,5 pt)$\\
\item 1) Déterminer le temps moyen consacré à la télévision le dimanche.
\item 2) Construire la courbe des fréquences cumulées croissantes encore appelé polygone
des fréquences cumulées croissantes\_\\
$(1,5 pt)$\\
\item 3) Déterminer la médiane de cette série statistique
pt)\\

\end{exercice}

% --- Fin Exercice inconnu ---


\newpage

% Exercice: Exercice inconnu
% Sujet: Géométrie
% Généré automatiquement

\begin{exercice}[Géométrie]
\label{ex:49}

Cxcicilc\\
(5 po$lnts)$\\
Lcs noles on Methomotiquos do 5^\{0\} oldvos €unu duuso do promloro D dun lycoo oblonuos opr^\{6\}s\\
qualriomg ovaluatlon onL^\{01\}oIooroupdos dong\\
lobloau suiunt\\
7^\{6\},2^\{4\}\\
Nolos\\
(0,'L\\
(8\\
[6,1^\{2\}/\\
\{1^\{2\},1^\{6\}\\
Elieculs\\
2^\{2\}\\
E:llectils cumules \_crolssants\\
\item a) Rocopicr ci comoloier 1otobloau predooonL
0,pt\\
Construire le polypono des plfoctlls aumulos coissants:\\
Delerminer la modiano do callo soi^\{1\} slatisuquu par Inlorpolation Inoaire. On prendra\\
cm\\
cmpour 5 unllos en ordonnoou.\\
0,7^\{5\} pi\\
pour 2 unilos cn obscissos 0l\\
Au momont do la croallon co lycoo 0^\{400\} éloves. On supposo quo chaque anneo\\
Dardo\\
7^\{5\}\% do s^\{08\} ancions éleves et qu^\{1\}l yp 3^\{00\}\\
apros sJ contion cot établlosomont\\
4^\{00\} @t pour lout entior nalurol non nul n, U, est l^\{0\}\\
nouvooux elovos: On nolo U,\\
nombro delovos do co lycoo aU cours do la n\\
annoo d'oxislenco\\
Cakulor U, ot U;\\
0,5 pt\\
Monbcr qu^\{0\} Un+\\
'Un + J^\{00\}.\\
0,5pi\\
1^\{200\}\\
On pose\\
Un pour lout onllor nalurel n non nul.\\
Domonber quo (Va) ost uno sulto oeomulrlque do raison\\
Oont on donncra l premior\\
0,7^\{8\} Pt\\
lermo\\
$Exprimor V. puis Un en lonctlon don.$\\
Ipt\\

\end{exercice}

% --- Fin Exercice inconnu ---


\newpage

% Exercice: Exercice inconnu
% Sujet: Mathématiques
% Généré automatiquement

\begin{exercice}[Mathématiques]
\label{ex:50}

PartleA :\\
5,5 points\\
Dans une classe de 5^\{0\} élèves, 1^\{5\} sont nés en 1^\{993\}, 2^\{0\} sont nés en\\
1^\{992\}, 1^\{0\} sont nés en 1^\{991\} et 5 sont nés en 1^\{990\}. On choisit au hasard et\\
simul$tanément 5 élèves de cette classe$\\
Tous les élèves ont la meme chance\\
dêtre choisis\\
Page\\
sur 2\\

\end{exercice}

% --- Fin Exercice inconnu ---


\newpage

% Exercice: Exercice inconnu
% Sujet: Mathématiques
% Généré automatiquement

\begin{exercice}[Mathématiques]
\label{ex:51}

Déterrniner\\
0,Spt\\
Le nombre de choix possibles\\
a)\\
0,7^\{5\}pt\\
Le nombre de choix où les 5 élèves ont le même age.\\
b)\\
0,7^\{5\}pt\\
Le nombre de choix où 4 élèves exactement ont le même\\
c)\\
âae\_\\

\end{exercice}

% --- Fin Exercice inconnu ---


\newpage

% Exercice: Exercice inconnu
% Sujet: Mathématiques
% Généré automatiquement

\begin{exercice}[Mathématiques]
\label{ex:52}

Les 5^\{0\} élèves ci-dessus sont candidats à un examen note sur\\
2-\\
\item 1^\{00\}. Après les résultats on obtient le tableau ci-dessous
Notes\\
[0 2^\{07\}\\
[2o,4^\{07\}\\
[4^\{0607\}\\
[6^\{0\},BoD\\
[8^\{0\}, 1OL\\
Effectfs\\
2^\{5\}\\
1^\{0\}\\

\end{exercice}

% --- Fin Exercice inconnu ---


\newpage

% Exercice: Exercice inconnu
% Sujet: Analyse
% Généré automatiquement

\begin{exercice}[Analyse]
\label{ex:53}

Construire le polygone des effectifs cumulés croissants de\\
cette série statistique .\\
1,Spt\\
Par lecture graphique, déterminer un encadrement\\
b-\\
damplitude 1^\{0\} de la note médiane de cette série.\\
0,Spt\\
Calculer à 1^\{0\}* près par défaut la moyenne et l'écart-type de\\
1,5pt\\
cette série\\

\end{exercice}

% --- Fin Exercice inconnu ---


\newpage

% Exercice: Exercice inconnu
% Sujet: Algèbre
% Généré automatiquement

\begin{exercice}[Algèbre]
\label{ex:54}

Partie B\\
x' -2x. (C) est la courbe représentative de\\
On considère la fonction 9 définie par g(x)\\
dans un repère orthonormé du plan.\\
Étudier les variations de 9 et tracer (C).\\
2,5\\
Déduire du tracer de (C) le tracé de la courbe de la fonction\\
Ig(x)\\
définie par\\
h(x)\\

\end{exercice}

% --- Fin Exercice inconnu ---


\newpage

% Exercice: Exercice inconnu
% Sujet: Algèbre
% Généré automatiquement

\begin{exercice}[Algèbre]
\label{ex:55}

$f(x) = -x^{3}$\\
Soit\\
la fonction définie sur IR par\\
(o;3).\\
On désigne par (C) sa courbe représentative dans un repère orthonormé\\
Que peut-on en déduire pour la courbe (C)\\
Etudier la parité de la fonction\\
Etudier les variations de\\
et dresser son tableau de variation.\\
Soit (T) la $tangente$\\
(C) en 0\\
Donner une équation de (T)\\
Etudier la position de (C) par rapport\\
(T)\\

\end{exercice}

% --- Fin Exercice inconnu ---


\newpage

% Exercice: Exercice inconnu
% Sujet: Algèbre
% Généré automatiquement

\begin{exercice}[Algèbre]
\label{ex:56}

parité\\
Etudier la\\
dérivée de\\
de la fonction\\
(C) aux points d'abscisses\\
Soit\\
un réel non nul. Montrer que les $tangentes$\\
respectives\\
sont parallèles\\
et -u\\
Tracer la courbe (C) et la $tangente (T) -$\\
(C) pour résoudre dans IR V'inéquation\\
0 < -x' + 3x < 2\\
Utiliser la courbe\\

\end{exercice}

% --- Fin Exercice inconnu ---


\newpage

% Exercice: Exercice inconnu
% Sujet: Analyse
% Généré automatiquement

\begin{exercice}[Analyse]
\label{ex:57}

Partie A\\
définie et dérivable sur R vérifiant les hypothèses suivantes\\
est une fonction impaire\\
La $limite en -c$\\
est +c\\
$f(2) = -2$\\
égal à\\
Le nombre dérivé de\\
en -1 est\\
Le nombre dérivé de\\
2 est egal\\
Le tableau incomplet des signes de la fonction dérivée de\\
est le suivant\\
+0\\
f'(x)\\

\end{exercice}

% --- Fin Exercice inconnu ---


\newpage

% Exercice: Exercice inconnu
% Sujet: Analyse
% Généré automatiquement

\begin{exercice}[Analyse]
\label{ex:58}

Recopier et compléter ce tableau des signes en met$tant le signe qui convient dans l'espace$\\
en pointillés\\
Déterminer f'(1), f'(2), f(0) et la $limite de$\\
1 pt\\
en +0\\
admet-elle un extremum relatif en -2\\
Justifier vos réponses\\
En 0\\
1 pt\\
pour tout nombre réel x, f(x)\\
ax;\\
On suppose que\\
bx\\
ou a, b et\\
sont des\\
coefficients reels.\\
Déterminer\\
1,5 pt\\

\end{exercice}

% --- Fin Exercice inconnu ---


\newpage

% Exercice: Probleme 11
% Sujet: Algèbre
% Généré automatiquement

\begin{exercice}[Algèbre]
\label{ex:59}

Problème\\
1^\{1\} points\\
Partle A\\
3,5 polnts\\
1^\{0\} 0^\{00\}\\
x?\\
définie par f(x)\\
est la fonction numérique dune variable réelle\\
7^\{3\})\\
On prend Icm\\
est sa courbe représentative dans le repère orthogonal (0,\\
pour 5^\{0\} sur l'axe des abscisses et Icm\\
0^\{00\} sur l'axe des ordonnées\\
pour\\
Etudier les variations de\\
sur [-1^\{00\}\\
1^\{00\}] et dresser son tableau de variations\_\\
2^\{0\} Construire la courbe Cf pour les abscisses appartenant\\
1^\{00\}]\\
Vintervalle [-1^\{00\}\\
définie par g(x)\\
f(x) + 1\\
3^\{0\} Soit\\
la fonction numérique de la variable réelle\\
$(0,7,3)$\\
Tracer la courbe Cg de g dans le repère\\
(On indiquera le programme de construction de C, à partir de Cf)\\

\end{exercice}

% --- Fin Probleme 11 ---


\newpage

% Exercice: Serie S
% Sujet: Géométrie
% Généré automatiquement

\begin{exercice}[Géométrie]
\label{ex:60}

1- Déterminer l'arrondi d'ordre 2 de la moyenne des dis$tances hebdomadaires$\\
0,7^\{5\} pt\\
parcourues par ces élèves\\
Déterminer la classe modale de cette série statistique.\\
0,5 pt\\
Dresser le tableau des effectifs cumules croissants\\
0,7^\{5\}pt\\
Déterminer par interpolation linéaire\\
la médiane de cette série statistique\\
nombre délèves qui parcourent moins de 5 km par semaine\\
Déterminer\\
0,5 pt\\
le professeur titulaire de\\
mieux\\
En vue de\\
cette classe voudrait constituer des groupes détude de cinq élèves,\\
Pour chacune des deux questions suivantes, quatre réponses sont proposées parmi\\
suivi\\
lesquelles une seule est juste. Ecrire\\
numéro de la question\\
lettre\\
sur votre feuille de composition\\
Aucune justification\\
juste\\
correspondant à la réponse\\
n'est demandée\\
Le nombre de groupes possibles qu'il peut former est\\
0 Spt\\
9^\{0\}-\\
\item b) A^\{3\}o
Le nombre de groupes qu'il peut former contenant au moins deux filles et au moins\\
deux garçons est\\
C^\{3\}o x C^\{3\}o\\
C^\{3\}o * C^\{3\}o\\
C^\{3\}o\\
C^\{3\}o x C^\{3\}o\\
Abo X A^\{3\}o\\
\item d) C^\{3\}o
Ipt\\

\end{exercice}

% --- Fin Serie S ---


\newpage

\end{document}
